% Part: normal-modal-logic
% Chapter: sequent-calculus
% Section: introduction

\documentclass[../../../include/open-logic-section]{subfiles}

\begin{document}

\olfileid[pt]{nml}{seq}{int}

\olsection{Introdução}

O cálculo de sequentes para a lógica proposicional pode ser estendido
por regras adicionais que tratam
de~\iftag{prvBox}{$\Box$\iftag{prvDiamond}{ e~}{}}{}%
\iftag{prvDiamond}{$\Diamond$}{}. Por exemplo, para~$\Log{K}$, temos
\Log{LK} acrescido de:
\[
  \iftag{prvBox}{\iftag{prvDiamond}{% <> and [] primitive
    \Axiom$\Gamma \fCenter \Delta, !A$
    \RightLabel{$\Box$}
    \UnaryInf$\Box\Gamma \fCenter \Diamond\Delta, \Box!A$
    \DisplayProof  
    \qquad
    \Axiom$!A, \Gamma \fCenter \Delta$
    \RightLabel{$\Diamond$}
    \UnaryInf$\Diamond!A, \Box\Gamma \fCenter \Diamond\Delta$
    \DisplayProof
}{% only Box primitive
\Axiom$\Gamma \fCenter !A$
\RightLabel{$\Box$}
\UnaryInf$\Box\Gamma \fCenter \Box!A$
\DisplayProof
}}{% only <> primitive
  \Axiom$!A \fCenter \Delta$
  \RightLabel{$\Diamond$}
  \UnaryInf$\Diamond!A \fCenter \Diamond\Delta$
  \DisplayProof
}
\]
Para extensões de~$\Log{K}$, também é preciso acrescentar regras
adicionais.

Nem toda lógica modal possui um cálculo de sequentes desse tipo. Até
mesmo $\Log{S5}$, que é semanticamente simples (pode ser definida sem
usar relações de acessibilidade), não é conhecida por possuir um
cálculo de sequentes que resulte de~$\Log{LK}$ e seja completo sem a
regra~\Cut. Contudo, ela possui um cálculo de \emph{hipersequentes}
completo e livre de corte.

\end{document}
